\section{Базовые принципы теории переговоров}

В теории игр процесс переговоров моделируется как задача о разделении некоторого ресурса (выигрыша) между участниками. Фундаментальной концепцией в анализе переговоров является определение статуса-кво и размера излишка, который создается при успешном исходе торгов.

\begin{definition}
\textbf{BANTA} (Best Alternative To a Negotiated Agreement) — лучшая альтернатива переговорному соглашению. Это полезность (выигрыш), которую игрок гарантированно получает в случае срыва переговоров. Также используются термины \textit{точка несогласия} или \textit{статус-кво}.
\end{definition}

Общий алгоритм решения переговорной задачи сводится к следующим шагам:
\begin{enumerate}
    \item Определение статуса-кво (BANTA) для каждого участника.
    \item Вычисление истинного размера «пирога» — суммарной выгоды, которая создается исключительно благодаря кооперации с учетом статуса-кво.
    \item Разделение размера «пирога» между участниками (базово — поровну, при допущении о равной переговорной силе игроков).
\end{enumerate}

Рассмотрим формальный пример совместного проекта. Пусть стоимость самостоятельной реализации проекта для первого игрока равна $X$, а для второго — $Y$. Совместная реализация проекта обходится в сумму $M$, причем $M < X + Y$. 
Размер «пирога» (общая экономия) составляет $X + Y - M$. При равном разделении каждый игрок получает экономию $\frac{X + Y - M}{2}$.

\begin{example}
Две компании (из Хьюстона и Сан-Франциско) пользуются услугами одного юриста из Нью-Йорка. Возможны два сценария оплаты перелетов:
\begin{itemize}
    \item \textit{Раздельные поездки}: Стоимость билетов Нью-Йорк — Хьюстон составляет $666\$$ в обе стороны ($2 \times 666 = 1332\$$). Стоимость билетов Нью-Йорк — Сан-Франциско составляет $1243\$$ в обе стороны ($2 \times 1243 = 2486\$$). Итоговые суммарные затраты: $1332 + 2486 = 3818\$$.
    \item \textit{Треугольный маршрут} (координация): Юрист летит по маршруту Нью-Йорк — Хьюстон — Сан-Франциско — Нью-Йорк. Общая стоимость составляет $2818\$$.
\end{itemize}

Определим выигрыш от координации (размер «пирога»): 
\[
3818 - 2818 = 1000\$
\]
При равном разделении «пирога» каждая фирма сэкономит по $500\$$. Следовательно, итоговые платежи составят:
\begin{itemize}
    \item Для фирмы из Хьюстона: $1332 - 500 = 832\$$
    \item Для фирмы из Сан-Франциско: $2486 - 500 = 1986\$$
\end{itemize}
\end{example}

\subsection{Стратегическое изменение статуса-кво}

Статус-кво имеет критическое значение не только сам по себе, но и в сравнении с альтернативой оппонента. В рамках стратегического взаимодействия игрокам может быть выгодно намеренно ухудшить собственную альтернативу, если это нанесет непропорционально больший урон альтернативе оппонента.

\begin{example}
Рассмотрим переговоры профсоюза и управляющего отелем. Изначально отель может заработать $1000\$$ в день. У работников есть опция заработать на стороне $300\$$ в день (их BANTA). Отель без работников зарабатывает $500\$$ (его BANTA). 
\begin{itemize}
    \item \textit{Исходный вариант}: Статус-кво $(300, 500)$. Размер «пирога» равен $1000 - 300 - 500 = 200\$$. При делении пирога поровну итоговый дележ составит $(300+100, 500+100) = (400, 600)$.
    \item \textit{Новый вариант}: Работники отказываются от части заработка на стороне, получая $200\$$, а освободившееся время тратят на забастовку у входа в отель. Это отпугивает клиентов, снижая BANTA отеля до $300\$$. Новый статус-кво $(200, 300)$. Размер «пирога» становится $1000 - 200 - 300 = 500\$$. Итоговый дележ: $(200+250, 300+250) = (450, 550)$.
\end{itemize}
Стратегия забастовки, несмотря на снижение собственной внепереговорной полезности, приводит к увеличению итогового выигрыша профсоюза.
\end{example}

\section{Модель торга с последовательными предложениями}

Классическая модель торга предполагает процесс, растянутый во времени, где участники делают предложения поочередно. Вводится фактор нетерпеливости (дисконтирования): со временем ценность одного и того же ресурса для игроков снижается.

Пусть торг длится $T$ периодов. Игроки имеют дисконт-факторы $\delta_1, \delta_2 \in (0, 1)$. Полезность получения доли $x$ в период $t$ для игрока $i$ составляет $x \delta_i^{t-1}$. Это означает, что получение всего ресурса после $t$ раундов переговоров эквивалентно получению доли $\delta_i^t$ немедленно.

\subsection{Конечный горизонт переговоров}

Рассмотрим игру с $T=2$. Игрок 1 делает предложение в первый период. Если Игрок 2 отказывается, он делает встречное предложение во второй период (игра «Ультиматум»).
Во втором периоде Игрок 2 предложит дележ $(0, 1)$, забирая всё себе. Его полезность, приведенная к началу игры, составит $\delta_2$, а полезность Игрока 1 будет $0$.
Следовательно, в первый период Игрок 1, применяя обратную индукцию, должен предложить Игроку 2 ровно его дисконтированную долю $\delta_2$, оставив себе $1 - \delta_2$.

В совершенном по подыграм равновесии (SPE) Игрок 1 предлагает дележ $(1 - \delta_2, \delta_2)$, и Игрок 2 его сразу принимает. 

\textbf{Выводы для конечной модели:}
\begin{itemize}
    \item Деление поровну не является обязательным.
    \item Чем терпеливее Игрок 2 (чем ближе $\delta_2$ к $1$), тем больше его равновесная доля.
    \item В двухпериодной модели терпеливость Игрока 1 не имеет значения, так как он не сталкивается с угрозой ожидания.
\end{itemize}

\subsection{Бесконечный горизонт переговоров}

Рассмотрим игру, где $T \to \infty$. Предлагающие чередуются бесконечно (в нечетные периоды предлагает Игрок 1, в четные — Игрок 2). В игре присутствуют два типа самоподобных подыгр. Решение ищется в классе стационарных стратегий: игроки всегда предлагают одно и то же распределение, когда наступает их очередь.

Пусть Игрок 1 всегда предлагает дележ $\mathbf{x}^* = (x_1^*, x_2^*)$ и принимает любое предложение $\mathbf{y} = (y_1, y_2)$, где $y_1 \ge y_1^*$. 
Пусть Игрок 2 всегда предлагает дележ $\mathbf{y}^* = (y_1^*, y_2^*)$ и принимает любое предложение $\mathbf{x} = (x_1, x_2)$, где $x_2 \ge x_2^*$.

\begin{theorem}
В модели бесконечного торга с чередующимися предложениями существует единственное совершенное по подыграм равновесие в стационарных стратегиях, в котором соглашение достигается в первый период с дележом:
\begin{equation}
    x_1^* = \frac{1 - \delta_2}{1 - \delta_1 \delta_2}, \quad x_2^* = \frac{\delta_2(1 - \delta_1)}{1 - \delta_1 \delta_2}
\end{equation}
\end{theorem}

\begin{proof}
Рассмотрим условия принятия предложений в соответствующих подыграх:
\begin{itemize}
    \item \textbf{Подыгра 1 (Предлагает Игрок 1):} Игрок 2 согласится, если предложенная доля не меньше его дисконтированной выгоды от продолжения игры: $x_2^* = \delta_2 y_2^*$.
    \item \textbf{Подыгра 2 (Предлагает Игрок 2):} Игрок 1 согласится, если $y_1^* = \delta_1 x_1^*$.
\end{itemize}
Учитывая балансовые ограничения $x_1^* + x_2^* = 1$ и $y_1^* + y_2^* = 1$, составим систему уравнений:
\begin{align}
    1 - x_1^* &= \delta_2 (1 - y_1^*) \\
    y_1^* &= \delta_1 x_1^*
\end{align}
Подставляя (2) в (1), получаем:
\[
    1 - x_1^* = \delta_2 (1 - \delta_1 x_1^*) \implies x_1^* (1 - \delta_1 \delta_2) = 1 - \delta_2 \implies x_1^* = \frac{1 - \delta_2}{1 - \delta_1 \delta_2}
\]
Доля второго игрока вычисляется как $x_2^* = 1 - x_1^* = \frac{\delta_2(1 - \delta_1)}{1 - \delta_1 \delta_2}$.
\end{proof}

\textbf{Свойства равновесия:}
\begin{itemize}
    \item \textit{Эффективность:} Согласие достигается немедленно (в первый период), ресурс не уничтожается из-за задержек.
    \item \textit{Преимущество первого хода:} Если игроки обладают одинаковой терпеливостью ($\delta_1 = \delta_2 = \delta$), то $x_1^* = \frac{1}{1 + \delta} > \frac{\delta}{1 + \delta} = x_2^*$.
    \item \textit{Эффект терпеливости:} Рост дисконт-фактора игрока (приближение к $1$) монотонно увеличивает его равновесную долю.
\end{itemize}

\begin{example}
Два игрока делят сумму $100\,000$. Дисконт-факторы равны $\delta_1 = 0.8$ и $\delta_2 = 0.9$. 
Обозначим дележи как $(x, 100\,000 - x)$ при предложении Игрока 1 и $(y, 100\,000 - y)$ при предложении Игрока 2.
Условия согласия:
\begin{align*}
    100\,000 - x &= 0.9(100\,000 - y) \\
    y &= 0.8x
\end{align*}
Подставляя $y$, получаем уравнение:
\[
    100\,000 - x = 90\,000 - 0.72x \implies 0.28x = 10\,000 \implies x \approx 35\,714
\]
Игрок 1 получит $35\,714$, Игрок 2 получит $64\,286$. В данном случае эффект высокой терпеливости второго игрока перевешивает преимущество первого хода Игрока 1.
\end{example}

\section{Переговоры с внешним риском остановки}

Альтернативной моделью является торг без дисконтирования во времени, но с экзогенным риском срыва. Пусть игроки делят сумму $M$. Если в текущем периоде соглашение не достигнуто, с вероятностью $p$ торг прекращается навсегда. В этом случае игроки получают свои статусы-кво $A$ и $B$ соответственно, причем $A + B < M$.

Составим условия индифферентности для принятия предложений. Игрок 1 предлагает дележ $(x, M-x)$, Игрок 2 предлагает $(y, M-y)$.
Игрок 2 согласится на предложение Игрока 1, если математическое ожидание от продолжения не превышает предложенную долю:
\begin{equation}
    M - x = p B + (1 - p)(M - y)
\end{equation}
Аналогично для Игрока 1:
\begin{equation}
    y = p A + (1 - p)x
\end{equation}
Подставляя $y$ из второго уравнения в первое, находим равновесную долю Игрока 1:
\begin{align*}
    M - x &= p B + (1 - p)\big(M - p A - (1 - p)x\big) \\
    x - (1 - p)^2 x &= M - p B - (1 - p)M + (1 - p)p A \\
    x \big(1 - (1 - p)^2\big) &= p M - p B + p A - p^2 A
\end{align*}
Разделив на $p$, получаем (приближенно для малых $p$ второстепенными членами пренебрегаем, либо в точном виде при строгом раскрытии):
\begin{equation}
    x = \frac{p M - p B + p A}{1 - (1 - p)^2} = \frac{M - B + A}{2 - p}
\end{equation}

Анализ предела при $p \to 0$:
\[
    \lim_{p \to 0} x = \frac{M - B + A}{2} = A + \frac{M - A - B}{2}
\]
При стремящейся к нулю вероятности срыва переговоров результат в точности совпадает с базовым принципом: каждый игрок получает свой статус-кво плюс ровно половину от созданного «пирога» $M - A - B$. При значимых значениях $p > 0$ возникает преимущество первого хода, так как риск потери ресурса заставляет второго игрока быть более сговорчивым.
